%! Modeling with Exponential Functions: Interest, Half-Life, and e — Unit 7, Lesson 7.4 — Guided Notes (Answer Key)
\documentclass[10pt]{article}
\usepackage{algebra2-article}
\usepackage{algebra2-key}   % pulls in -boxes; do NOT also load -boxes
\newcommand{\TallMath}[1]{$\displaystyle #1\rule[-1.4em]{0pt}{3.2em}$}
% Growth curve on a money-sized grid. pgfmath's pow() is unreliable for fractional
% exponents, so every exponential is plotted as a*exp((ln b)t).
% ln(1.05) = 0.0487902   ln(1.06) = 0.0582689   ln(0.5) = -0.693147   ln(2) = 0.693147
% Vocab answer for the key: bold forest term, then the filled definition on a write-line.
% The leading and trailing \par are required: \ansline ends with \dotfill but does NOT end the
% paragraph, so without them each term label gets pulled onto the previous answer's line.
\newcommand{\vocabans}[2]{%
  \par\noindent\textbf{\textcolor{forest}{#1:}}\\[1pt]\ansline{#2}\par}

\begin{document}

\pageheader{Unit 7, Lesson 7.4}{Guided Notes --- Answer Key}
\namedateperiod

\vspace{0.08in}

\begin{objectivebox}
\textbf{By the end of this lesson, I will be able to\dots}
\begin{itemize}\setlength{\itemsep}{2pt}
  \item build and evaluate an exponential model for compound interest, continuous growth, half-life, and doubling --- and say what the answer \emph{means} in the situation;
  \item explain why the base is always the factor for \emph{one period} and the exponent always \emph{counts periods}, so the rate is divided and the time is multiplied;
  \item describe where the number $e$ comes from and use $A=Pe^{rt}$;
  \item recognize an exponential equation that no common base will ever solve, and locate its answer on a graph anyway.
\end{itemize}
\end{objectivebox}

\vspace{0.05in}

\begin{vocabbox}
\small
Fill in each term as we build it together.
\par\vspace{2pt}   % \par is required: \vocabans's \noindent is a no-op mid-paragraph
\vocabans{Principal, $P$}{the starting amount --- the money deposited or borrowed before any interest; it is the $a$ of Lesson 7.1's $y=ab^{x}$}
\vocabans{Compounding period}{the stretch of time between two moments when interest is added; $n$ of them fit in a year, so each one uses the rate $\frac{r}{n}$}
\vocabans{Compound interest formula, $A=P\left(1+\frac{r}{n}\right)^{nt}$}{the base is the factor for one period and the exponent counts the periods in $t$ years}
\vocabans{The number $e$, and continuous compounding $A=Pe^{rt}$}{$e\approx2.71828$ is the value $\left(1+\frac1m\right)^{m}$ closes in on as $m$ grows; $Pe^{rt}$ is the ceiling --- compounding chopped infinitely fine}
\vocabans{Half-life}{the time it takes for half of a quantity to be gone; if it is $h$, the model is $A_{0}\left(\frac12\right)^{t/h}$}
\vocabans{Doubling time}{the time it takes for a quantity to double; if it is $d$, the model is $A_{0}\,2^{\,t/d}$}
\end{vocabbox}

\begin{hookbox}
Four banks, one deposit of \$1000, one year, and the same advertised rate of $6\%$. The only difference
is how often each bank stops and adds the interest in.

\vspace{4pt}
\begin{center}
\small
\renewcommand{\arraystretch}{1.35}
\begin{tabular}{|l|c|c|}
\hline
\textbf{Bank} & \textbf{Interest added} & \textbf{Balance after 1 year} \\
\hline
Ashmore     & once a year (annually) & \$1060.00 \\
Bellhaven   & 4 times a year (quarterly) & \$1061.36 \\
Corliss     & 12 times a year (monthly) & \$1061.68 \\
Dunmore     & 365 times a year (daily) & \$1061.83 \\
\hline
\end{tabular}
\end{center}

\vspace{4pt}
\begin{itemize}\setlength{\itemsep}{2pt}
  \item Dunmore adds interest \textbf{365 times} as often as Ashmore. How much more does that earn you over the year? \ans{\$1.83}
  \item Going from once a year to four times a year gained \ans{\$1.36}. Going from 12 times to 365 times gained \ans{\$0.15}. What is happening to the size of the gains? \ansline{They are shrinking fast --- each extra chop of the year buys less than the one before it.}
  \item Predict: is there a bank so aggressive --- compounding every second, every instant --- that it doubles your money in a year? Or is something stopping it? \ansline{Something is stopping it: the balances are closing in on a ceiling near \$1061.84, not running away. (Answers will vary; the point is the prediction.)}
\end{itemize}
\end{hookbox}

\boxguard
\begin{notesbox}{1. Cutting the Year Into Pieces}
Lesson 7.1 built $y=ab^{x}$, where $b$ was the growth factor for \emph{one} year and $x$ counted
\emph{years}. Nothing about that changes today --- except that the ``one period'' does not have to be a
year.

\vspace{4pt}
If an account pays an annual rate $r$ but adds the interest $n$ times per year, then each period uses
only a fraction of the rate, and there are more periods to count:

\vspace{2pt}
\begin{itemize}[leftmargin=1.4em, itemsep=2pt]
  \item the rate for one period is \ans{$\frac{r}{n}$}, so the \textbf{factor} for one period is $1+$ \ans{$\frac{r}{n}$};
  \item in $t$ years the number of periods is \ans{$nt$}.
\end{itemize}

\vspace{4pt}
\begin{center}
\begin{tcolorbox}[colback=goldbg, colframe=goldacc, boxrule=0.8pt, arc=1.5mm, width=0.92\linewidth,
    left=3mm, right=3mm, top=2mm, bottom=2mm]
\centering
\textbf{Compound Interest} \\[3pt]
\TallMath{A = P\left(1+\frac{r}{n}\right)^{nt}}
\\[1pt]
{\small $P$ = principal (starting amount) \quad $r$ = annual rate as a \emph{decimal} \quad
$n$ = compoundings per year \quad $t$ = years}
\end{tcolorbox}
\end{center}

\vspace{2pt}
\textbf{Read the formula as one sentence:} the base $\left(1+\frac{r}{n}\right)$ is what happens in
\ans{one period}, and the exponent $nt$ \ans{counts the periods}. Every error today is a place where those two stop
agreeing with each other.

\vspace{5pt}
\textbf{Worked example.} \$2500 is invested at $4\%$ compounded quarterly for $6$ years.

\vspace{3pt}
\begin{center}
\small
\renewcommand{\arraystretch}{1.5}
\begin{tabularx}{0.95\linewidth}{@{}l X@{}}
\hline
$P=$ \ans{2500} \quad $r=$ \ans{0.04} \quad $n=$ \ans{4} \quad $t=$ \ans{6}
  & \emph{$4\%$ becomes $0.04$ --- Lesson 7.1's translation, unchanged.} \\
Rate per quarter: $\frac{r}{n}=$ \ans{0.01} \quad Base: \ans{1.01}
  & \emph{One quarter's worth of interest, not one year's.} \\
Number of quarters: $nt=$ \ans{24}
  & \emph{Six years of quarters --- the exponent counts periods.} \\
$A=$ \ans{$2500(1.01)^{24}$} $=$ \ans{\$3174.34}
  & \emph{Round money to the nearest cent, only at the end.} \\
\hline
\end{tabularx}
\end{center}

\vspace{2pt}
The same \$2500 at $4\%$ compounded \emph{annually} for 6 years would be
$2500(1.04)^{6}=\$3163.30$. Compounding quarterly is worth an extra \ans{\$11.04} --- real, but small.
\end{notesbox}

\boxguard
\begin{notesbox}{2. Where the Gains Stop: the Number $e$}
The Hook found gains that shrink fast. To see what they shrink \emph{toward}, strip the money away and
ask the cleanest version of the question: start with \$1 at $100\%$ for one year, and chop the year into
$m$ equal pieces. The balance is $\left(1+\frac1m\right)^{m}$.

\vspace{3pt}
\begin{center}
\small
\renewcommand{\arraystretch}{1.4}
\begin{tabular}{|l|c|c|c|c|c|c|}
\hline
\textbf{pieces $m$} & $1$ & $10$ & $100$ & $1000$ & $10{,}000$ & $100{,}000$ \\
\hline
$\left(1+\frac1m\right)^{m}$ & $2$ & $2.59374$ & $2.70481$ & \ans{2.71692} & \ans{2.71815} & \ans{2.71827} \\
\hline
\end{tabular}
\end{center}

\vspace{3pt}
The values climb, but they are not climbing \emph{toward infinity} --- they are closing in on one
number. That number is called \ans{$e$}, and to five decimal places it is \ans{2.71828}. Like
$\pi$, it never terminates and never repeats, and like $\pi$, it has a button on your calculator.

\vspace{4pt}
\begin{center}
\begin{tcolorbox}[colback=goldbg, colframe=goldacc, boxrule=0.8pt, arc=1.5mm, width=0.92\linewidth,
    left=3mm, right=3mm, top=2mm, bottom=2mm]
\centering
\textbf{Continuous Compounding} \\[3pt]
\TallMath{A = Pe^{\,rt}}
\\[1pt]
{\small No $n$ appears --- the compounding has been chopped so fine that counting periods is over.}
\end{tcolorbox}
\end{center}

\vspace{2pt}
\textbf{Same investment, no more chopping.} \$2500 at $4\%$ compounded \emph{continuously} for 6 years:
$A=2500e^{\,0.04\cdot6}=2500e^{\,0.24}=$ \ans{\$3178.12}.

\vspace{3pt}
Compare the three answers for the identical deal: annually \$3163.30, quarterly \$3174.34, continuously
\ans{\$3178.12}. Continuous compounding is the \textbf{ceiling} --- the most that rate can possibly
produce --- and it beats quarterly by only \ans{\$3.78}. This answers the Hook: no bank can double
your money by compounding more often, because $e^{\,r}$ is where the chopping runs out.

\vspace{2pt}
\emph{Why anyone cares:} populations, cooling, and radioactive decay do not wait for the end of a
quarter. They change continuously, so $e$ is their natural base --- which is why it comes back in Unit 8.
\end{notesbox}

\boxguard
\begin{notesbox}{3. Half-Life and Doubling Time: the Same Idea on a Different Clock}
A \textbf{half-life} is the time it takes for half of a quantity to disappear. A \textbf{doubling time}
is the time it takes for a quantity to double. Both name \emph{how long one period lasts}, and both
models are built exactly the way the interest formula was.

\vspace{4pt}
\begin{center}
\begin{tcolorbox}[colback=sky, colframe=navy, boxrule=0.8pt, arc=1.5mm, width=0.94\linewidth,
    left=3mm, right=3mm, top=2mm, bottom=2mm]
\centering
\textbf{Half-life $h$:} \; \TallMath{A(t) = A_{0}\left(\tfrac12\right)^{t/h}}
\qquad
\textbf{Doubling time $d$:} \; \TallMath{A(t) = A_{0}\,2^{\,t/d}}
\\[1pt]
{\small The base says what one period does. The exponent $t/h$ (or $t/d$) counts how many periods have
gone by.}
\end{tcolorbox}
\end{center}

\vspace{2pt}
\textbf{Why $t/h$ and not $t$?} Because $t$ is measured in hours or years, and the exponent must be
measured in \emph{periods}. If the half-life is 6 hours, then in 18 hours there have been
\ans{3} halvings, and in 9 hours there have been \ans{1.5} --- the Warm-Up answer. Dividing
by $h$ is the unit conversion.

\vspace{5pt}
\textbf{Worked example.} A hospital gives a patient 80 mg of a tracer whose half-life is 6 hours.

\vspace{3pt}
\hspace*{1.0em} Model: $A(t)=$ \ans{$80\left(\frac12\right)^{t/6}$} \quad where $t$ is in \ans{hours}.

\vspace{4pt}
\begin{center}
\small
\renewcommand{\arraystretch}{1.4}
\begin{tabular}{|l|c|c|c|c|c|}
\hline
$t$ (hours) & $0$ & $6$ & $12$ & $18$ & $9$ \\ \hline
halvings $t/6$ & $0$ & \ans{1} & \ans{2} & \ans{3} & \ans{1.5} \\ \hline
$A(t)$ (mg) & $80$ & \ans{40} & \ans{20} & \ans{10} & \ans{28.28} \\ \hline
\end{tabular}
\end{center}

\vspace{2pt}
The last column is the one worth arguing about: $t=9$ is \emph{between} two halvings, the exponent is
$1.5$, and the answer lands sensibly between 40 and 20.

\vspace{3pt}
\textbf{Will the tracer ever be completely gone?} \ans{No} --- and the reason is Lesson 7.0's, not
a new one: the range of this model is \ans{$(0,80]$} and the graph has a horizontal asymptote at
\ans{$A=0$}. Half of something is never nothing. \emph{(In context, of course, the amount eventually
falls below anything a machine can detect --- but the model never reaches zero.)}

\vspace{4pt}
\begin{tcolorbox}[colback=greenbg, colframe=greenacc, boxrule=0.8pt, arc=1.5mm,
    left=2.5mm, right=2.5mm, top=1.5mm, bottom=1.5mm]
\textbf{Yesterday's method still works here.} When does the tracer drop to 10 mg? Divide first, then
use a common base (Lesson 7.3).

\vspace{3pt}
\hspace*{1.0em} From $80\left(\frac12\right)^{t/6}=10$: \quad
$\left(\frac12\right)^{t/6}=$ \ans{$\frac18$} $=\left(\frac12\right)^{\blacksquare}$ with
$\blacksquare=$ \ans{3}, \quad so $\frac{t}{6}=$ \ans{3} \; and \; $t=$ \ans{18} hours.

\vspace{3pt}
That worked because $10$ is $80$ cut in half three times. Ask for 25 mg instead and the method dies.
\end{tcolorbox}
\end{notesbox}

\boxguard[30]
\begin{notesbox}{4. The Wall}
Here is the question every investor actually asks: \textbf{when does my money double?} Put \$1000 into
an account paying $5\%$ compounded annually. The model is $A(t)=1000(1.05)^{t}$, and doubling means
solving $1000(1.05)^{t}=2000$. Divide both sides by $1000$: \; $(1.05)^{t}=$ \ans{2}.

\vspace{3pt}

\begin{center}
\begin{minipage}[c]{0.46\linewidth}
\centering
\begin{tikzpicture}[x=0.26cm, y=0.00155cm]
  \draw[linegray!40, very thin, xstep=2, ystep=500] (0,0) grid (20,3000);
  \draw[->, semithick, charcoal] (0,0)--(21.5,0) node[above right, font=\scriptsize]{$t$ (yr)};
  \draw[->, semithick, charcoal] (0,0)--(0,3250) node[left, font=\scriptsize]{$A$};
  \foreach \t in {5,10,15,20} \node[below, font=\tiny] at (\t,0) {$\t$};
  \foreach \a in {1000,2000,3000} \node[left, font=\tiny] at (0,\a) {$\a$};
  \draw[very thick, forest, domain=0:20, samples=120, smooth]
    plot (\x,{1000*exp(0.0487902*(\x))});
  \draw[goldacc, very thick] (0,2000)--(20,2000);
  \node[goldacc, font=\scriptsize, above] at (5,2020) {$A=2000$};
  \fill[charcoal] (14,1979.9) circle (2.2pt);
  \fill[charcoal] (15,2078.9) circle (2.2pt);
\end{tikzpicture}
\end{minipage}
\begin{minipage}[c]{0.50\linewidth}
\small
\begin{itemize}[leftmargin=1.3em, itemsep=2pt]
  \item The gold line \emph{does} cross the curve, so the answer \textbf{exists}.
  \item $1000(1.05)^{14}=\$1979.93$ and $1000(1.05)^{15}=\$2078.93$, so the money doubles between year
        \ans{14} and year \ans{15}.
  \item Now try Step 1 of Lesson 7.3 on $(1.05)^{t}=2$: rewrite both sides over one base.
        \ans{$2$ is not a power of $1.05$, so there is no common base --- Step 1 is impossible.}
  \item So this is \emph{not} a ``no solution'' equation. Compare $2^{x}=-3$: that one has no solution
        at all, because $-3$ is not in the \ans{range}.
\end{itemize}
\end{minipage}
\end{center}

\vspace{3pt}
\begin{tcolorbox}[colback=redbg, colframe=redacc, boxrule=0.8pt, arc=1.5mm,
    left=2.5mm, right=2.5mm, top=1.5mm, bottom=1.5mm]
\textbf{The running list.} Every equation below has an answer. Not one of them has a common base.
\[
  2^{x}=5 \qquad 132(0.85)^{t}=32 \qquad (0.9)^{t}=2.118 \qquad (1.05)^{t}=2
\]
Half-life questions often \emph{do} land on a common base --- halving to $\frac18$ is what the base
does. Interest questions almost never do. The tool that finishes all four is a \textbf{logarithm}, and
it is the whole of Unit 8.
\end{tcolorbox}
\end{notesbox}

\begin{practicebox}
Show the substitution before you reach for the calculator. Round money to the nearest cent.
\begin{enumerate}[leftmargin=1.7em, itemsep=5pt]
  \item \$3000 is invested at $5\%$ compounded monthly for 4 years. \\[2pt]
        $\dfrac{r}{n}=$ \ans{$\frac{0.05}{12}$} \quad $nt=$ \ans{48} \quad
        $A=$ \ans{$3000\left(1+\frac{0.05}{12}\right)^{48}$} $=$ \ans{\$3662.69}

  \item The same \$3000 at $5\%$ compounded \emph{continuously} for 4 years. \\[2pt]
        $A=$ \ans{$3000e^{\,0.20}$} $=$ \ans{\$3664.21} \quad Extra earned by compounding continuously:
        \ans{\$1.52}

  \item A 240 mg dose of a drug has a half-life of 8 hours. Write the model, then find how much
        remains after 24 hours. \\[2pt]
        $A(t)=$ \ans{$240\left(\frac12\right)^{t/8}$} \hfill $A(24)=$ \ans{30 mg}

  \item Which of these can be solved with Lesson 7.3's common-base method, and which cannot? Circle
        one on each line and say how you can tell at a glance.
        \\[2pt]
        \hspace*{1.0em} $240\left(\frac12\right)^{t/8}=30$ \quad \textcolor{keyred}{\textbf{yes}} / \textbf{no}
        \hfill
        $5000(1.03)^{t}=10{,}000$ \quad \textbf{yes} / \textcolor{keyred}{\textbf{no}}
        \\[2pt]
        \ansline{Divide first: $\left(\frac12\right)^{t/8}=\frac18$ is a power of $\frac12$, but $(1.03)^{t}=2$ asks for $2$ as a power of $1.03$, and it is not one.}
\end{enumerate}
\end{practicebox}

\begin{teachernote}
\textbf{The Hook's numbers are the argument; do not compute them at the board.} They are printed
because the surprise lives in the \emph{differences}, not the arithmetic: \$1.36, then \$0.32, then
\$0.15. Students almost universally predict that daily compounding is dramatically better, and the
\$1.83 total gain over an entire year is what makes Section 2 feel necessary rather than decorative.
Ask the third bullet as a real prediction and write the guesses down --- Section 2 settles it.

\textbf{One sentence carries the whole lesson:} \emph{the base is what happens in one period, and the
exponent counts the periods.} Say it at the compound-interest box, again at the half-life box, and a
third time in Section 4. Both signature errors are that sentence broken in half --- using $r$ where
$\frac{r}{n}$ belongs (base kept annual, exponent made quarterly) or $t$ where $nt$ belongs (the
reverse). They come back as the two named students in the Activity, so name them here.

\textbf{Section 2's table is where $e$ has to be earned, not announced.} Let students see $2.71692$,
$2.71815$, $2.71827$ and say out loud that the digits are freezing from the left. Resist defining $e$
by a formula --- ``the number the chopping runs into'' is the honest description at this level, and it
is exactly what makes $A=Pe^{rt}$ read as a \emph{ceiling}. The $\$3.78$ gap between quarterly and
continuous is the punchline: infinitely fine compounding is worth less than four dollars.

\textbf{Section 3's $t=9$ column is the one to slow down on.} A fractional number of half-lives reads
as an error to most students --- the same instinct that made them distrust $x=\frac23$ in Lesson 7.3.
Point at the Warm-Up: they already computed $\left(\frac12\right)^{1.5}\approx0.354$, and
$28.28$ sits exactly where it should, between $40$ and $20$. The ``will it ever be gone?'' question is
pure Lesson 7.0 and should be answered from the asymptote, not from arithmetic.

\textbf{Section 4 must end unresolved, and students must not leave calling it ``no solution.''} The
graph does the work: the gold line crosses, so an answer exists; $2$ is simply not a power of $1.05$.
The contrast with $2^{x}=-3$ is the distinction Lesson 7.3's exit ticket was built to test, and it is
the one Unit 8 depends on. Do not name logarithms. Leave the four-equation list posted --- Lesson 8.4
takes it down.
\end{teachernote}

\end{document}
